\section*{الفصل \ref{ch:g8:plastics} --- البلاستيك والمواد الاصطناعية}

\begin{solution}{exo:g8:plastics:1}
طبيعية: القطن، والصوف. صناعية: الورق، والفسكوز. اصطناعية: النايلون،
والبوليستيرين.
\end{solution}

\begin{solution}{exo:g8:plastics:2}
من النفط الخام (أو الغاز الطبيعي).
\end{solution}

\begin{solution}{exo:g8:plastics:3}
PET (بولي إيثيلين تيريفثالات): قوارير الماء والمشروبات الغازية.
\end{solution}

\begin{solution}{exo:g8:plastics:4}
\omterm{def:g8:plastics:thermoplastic}{البلاستيك الحراري} يلين كلما سُخّن ويمكن إعادة قولبته؛ \omterm{def:g8:plastics:thermoplastic}{والبلاستيك المتصلّد
بالحرارة} يتصلب نهائيًا في المرة الأولى ويتفحم بدلًا من أن يلين.
\end{solution}

\begin{solution}{exo:g8:plastics:5}
البولي بروبيلين (PP)، وهو \omterm{def:g8:plastics:thermoplastic}{بلاستيك حراري}: يمكن صهره وإعادة قولبته.
\end{solution}

\begin{solution}{exo:g8:plastics:6}
المثلث يسمّي \omterm{def:g8:plastics:plastic}{البلاستيك} فقط، ليساعد على الفرز. أما إعادة تدوير الجسم فتتوقف
على جمعه وعلى وجود مصنع يعيد تدوير ذلك
\omterm{def:g8:plastics:plastic}{البلاستيك}.
\end{solution}

\begin{solution}{exo:g8:plastics:7}
المقلاة تسخن: \omterm{def:g8:plastics:thermoplastic}{فالبلاستيك الحراري} يلين، أما المتصلّد بالحرارة فلا.
\end{solution}

\begin{solution}{exo:g8:plastics:8}
$\frac{353}{460} \approx 0.77$: أي نحو \qty{77}{\percent}.
\end{solution}

\begin{solution}{exo:g8:plastics:9}
يحوي PVC الكلور: فحرقه يطلق كلوريد الهيدروجين، وهو غاز أكّال، إضافة إلى
السخام وأول أكسيد الكربون الناتجين عن \omterm{def:g8:combustion-and-fuels:complete}{الاحتراق غير التام} في نار
الحديقة.
\end{solution}

\begin{solution}{exo:g8:plastics:10}
$300 \times 36 = 10\,800$ قارورة؛ و$10\,800 \times 25 = 270\,000$\,g، أي
\qty{270}{kg}.
\end{solution}

\begin{solution}{exo:g8:plastics:11}
ثنائي أكسيد الكربون والماء. فكربونه يأتي من النفط الخام، المخزّن تحت الأرض
ملايين السنين: فحرقه يضيف ثنائي أكسيد الكربون إلى \omterm{def:g4:air-a-mixture-of-gases:air}{الهواء}، كما يفعل حرق
\omterm{def:g8:combustion-and-fuels:fossil}{الوقود الأحفوري}.
\end{solution}

\begin{solution}{exo:g8:plastics:12}
القارورة لا تتحلل: فالشمس والأمواج تكسّرها قطعًا أصغر فأصغر. وتطفو القطع
الدقيقة في الماء أو تغوص، فتبتلعها الأسماك مع
غذائها.
\end{solution}

\begin{solution}{pb:g8:plastics:1}
\textbf{1.} أجسام القوارير: 1 (PET). الأغطية: 2 (HDPE) أو 5 (PP).

\textbf{2.} اصطناعية: فهي مصنوعة من مواد بناها الكيميائيون، وفي الغالب من
النفط.

\textbf{3.} نعم: فمادة PET \omterm{def:g8:plastics:thermoplastic}{بلاستيك حراري}.

\textbf{4.} يُعاد تدوير كل \omterm{def:g8:plastics:plastic}{بلاستيك} على حدة: \omterm{def:g6:pure-substances-and-mixtures:mixture}{فخليط} من أنواع \omterm{def:g8:plastics:plastic}{البلاستيك}، إذا
صُهر معًا، يعطي مادة رديئة.

\textbf{5.} $0.80 \times 1000 = \qty{800}{kg}$.

\textbf{6.} الأغطية: $0.12 \times 1000 = \qty{120}{kg}$. الملصقات: $0.05
\times 1000 = \qty{50}{kg}$.

\textbf{7.} $0.03 \times 1000 = \qty{30}{kg}$ من الأوساخ والسوائل.

\textbf{8.} $80 + 12 + 5 = \qty{97}{\percent}$.

\textbf{9.} $0.025 \times 800 = \qty{20}{kg}$.

\textbf{10.} الرقائق توفر النفط الخام وطاقة صنع PET جديد، وتحوّل النفايات إلى
جسم نافع.

\textbf{11.} ثنائي أكسيد الكربون (وإذا كان \omterm{def:g4:what-a-fire-needs:burning}{الاحتراق} رديئًا، أول أكسيد الكربون
والسخام).

\textbf{12.} $800 - 20 = \qty{780}{kg}$ من رقائق PET النظيفة لكل رزمة.
\end{solution}
